The thesis set out to design a calibrated and realistic Agent-Based Model for central clearing, that includes a client tier, the workings of which was highlighted as an open question by \cite{menkveld_vuillemey_2021}. It adapts and extends the work of \cite{deloitte_ccp} and \cite{simudyne_odd} who built a simpler model to analyse CCP systemic risk. To this end the model was built up in two layers, a market simulation where traders following different strategies put in order onto the LOB of an exchange, and a clearing layer where a CCP novates these trades by these traders who either hold CCP membership or act as the clients of BCMs and NBCMs that do. The clearing layer consists of a complete margin framework - including variation margin, initial margin, and default fund contributions - and a default management framework mimicking a CCP's default waterfall. The model was calibrated to reproduce the stylised facts of returns identified by \cite{cont_2001}, based on E-mini S\&P future data during a calm and stressed period, reproducing moments such as heavy tails and volatility clustering. A surrogate modelling approach used for ABM calibration is employed, initially built by \cite{lamperti_2018} and adapted for financial market simulations by \cite{gao_2022_xgb}. The model itself serves as the thesis's primary contribution, and as a novel contribution to the best of our knowledge.

The model produces no defaults in the calm period, and a low amount of defaults in the stressed period. These are mostly client defaults, absorbed by their CM and never reaching the CCP. Looking at the output from this simplified and stylised model, a few observations can be drawn. In the stressed period, there is the occasional case of a NBCM default, caused by an inability to liquidate its own position to raise the funding required to meet margin calls. This highlights the importance for these NBCMs to have a readily available external source of funding which is robust to stressful market environments. Another observation is the tendency for margin calls to cluster in time, leading to short bursts during stressful periods. 

The model is then applied to study the effects of various kinds of flat margin levels in contrast to VaR style volatility reactive margining. Here, we see that high margining tends to relocate risk rather than remove it. As flat margin level rises, client defaults fall but member defaults rise. This is caused by capital ratio floors being reached earlier on, causing the members to freeze the provision of clearing services. Some clients are then trapped holding on to large positions that, upon client default, the member is no longer able to absorb. This highlights the importance for clients to have an alternative clearing connection that can be used in case their primary connection freezes their provision of clearing services. This is even more important considering the results of \cite{ofr_2026} who find the modal client to have only a single clearing connection and of \cite{acosta_smith_2018} who found that the introduction of capital and leverage ratio requirements led to a reduction in member willingness to provide clearing services. 

Finally, an alternative version of the model is presented that allows for the model to be run across a diverse set of stochastic synthetic paths. This allows for an expansion of the capabilities of the model to analyse CCP and member risks across a large set of scenarios. The calibration of this model was also a novel application of the methodologies of \cite{lamperti_2018} and \cite{gao_2022_xgb}. It also reproduces stylised facts to a comparable level as the standard model. The use of this synthetic alternative opens up a lot of avenues for further research. Parameters of the clearing layer can be systematically varied to study the effect of each change, the design of the clearing layer can be enriched to trial the introduction of new regulation, and it can be calibrated on other data sources and adapted to the regulatory guidelines of different asset types. This leaves an excellent avenue for future work. 

The model does have its own limitations. It is created as a realistic replication of central clearing, but is after all a stylised model. The strength of the model is then also reliant in the strength of the assumptions made during its construction. A second limitation is that the calibration was performed with the clearing layer inactive, as it is a computationally expensive step. The trading freezes and deleveraging of the clearing layer can, however, impact the market layer. A further limitation is that the fundamental price path is exogenous, and with the simulated mid tracking it closely, the effects of the clearing layer on prices only feed into transient transaction costs. The crash amplification channel through margin call driven fire sales is largely absent. Any observation drawn from the model is a stylised result, and cannot be attributed to causation like it can for standard data analysis, but the use of an ABM brings with it it's own benefits. Furthermore, the effects seen in the model are difficult to distinguish from the elements of model construction -- i.e. do we see this because it may be so in real markets or because the model was designed in this certain way. 




